\section{The previously answered question about both odd labels}
The final paragraph of the bound source asks whether, for every prime $p\ge5$, there is a regular polyhedron of order $2^n p$ with both adjacent orders divisible by $p$. This existence assertion is already supplied by the polygon construction of Montero--Ivi\'c Weiss and the later equal-parameter construction of Li--Zhang. We include the elementary group model to make this use independent of an unpaid second umbrella. We do not claim a new solution of that question, or incorporate the separate higher-rank conjecture.

Let $r\ge3$, put $\Omega=\mathbb Z/r\mathbb Z$, and let
\[
 V=\{(v_i)\in\mathbb F_2^\Omega:\sum_i v_i=0\}.
\]
The dihedral group $D=\langle x,b\rangle$ of order $2r$ acts by $x(i)=i+1$ and $b(i)=-i$. Put $s=xb$, $q=e_0+e_1$, and use the left multiplication
\[
 (v,g)(w,h)=(v+gw,gh)
\]
in $V\rtimes D$. The elements
\[
 \rho_0=(0,s),\qquad \rho_1=(0,b),\qquad \rho_2=(q,s)
\]
are distinct nonidentity involutions; the endpoints commute because $sq=q$. Their first adjacent product is $(0,x)$ of order $r$.

For $g\in D$, the vector $c(g)=e_0+ge_0$ lies in $V$ and satisfies $c(gh)=c(g)+gc(h)$. Therefore
\[
 \delta(v,g)=(v+c(g),g)
\]
is an automorphism, of square one, fixing $\rho_1$ and exchanging $\rho_0,\rho_2$. The second adjacent product consequently also has order $r$. Moreover
\[
 \langle\rho_0,\rho_1\rangle=\{(0,g):g\in D\},\qquad
 \langle\rho_1,\rho_2\rangle=\{(c(g),g):g\in D\}.
\]
Their intersection corresponds to $c(g)=0$, equivalently $g(0)=0$, whose stabilizer is $\langle b\rangle$. The rank-three test in [U] proves every marked subgroup intersection.

The product $\rho_0\rho_2$ is translation by $q$. Its conjugates under the first dihedral subgroup give every $e_i+e_{i+1}$. Summing these along a path from $0$ to $i$ gives $e_0+e_i$; these $r-1$ vectors form a basis of $V$. Hence the three markers generate the full semidirect product, of order $2^{r-1}\cdot2r=r2^r$. The standard regular-polytope correspondence and $\delta$ give a self-dual regular polyhedron of type $\{r,r\}$. Taking $r=p$ gives the requested order $p2^p$ and both adjacent labels $p$.

This proof reuses the existing equal-label component at precisely its existence interface. The full universal presentation and identification with the geometric halving operation are proved in the separately accepted Topic C application; they are not premises of the argument just given. Their additional source claims and costs are not silently imported or credited here. The present complete model, intersection and generation proof is charged within this application.
